\section*{Hoofdstuk \ref{ch:b1:elementary-steps} --- Reactiemechanismen: elementaire stappen en benaderingen}

\begin{solution}{exo:b1:elementary-steps:1}
(a) Unimoleculair: kan elementair zijn. (b) Bimoleculair: kan elementair
zijn. (c) Drie moleculen en vier bindingen die tegelijk worden
omgebouwd: niet elementair. (d) Trimoleculair, waarbij \ce{N2} de
vrijkomende energie afvoert: een zeldzame maar echte \omterm{def:b1:elementary-steps:elementary-step}{elementaire stap}.
\end{solution}

\begin{solution}{exo:b1:elementary-steps:2}
$v = k[\mathrm{A}][\mathrm{B}]$; $v = k[\mathrm{A}]^2$; $v = k[\mathrm{A}]$.
Een totale reactievergelijking is een balans, geen beschrijving van de
manier waarop de moleculen elkaar ontmoeten; haar \omterm{def:b1:rate-laws:order}{snelheidswet} moet
gemeten worden, of afgeleid uit een mechanisme.
\end{solution}

\begin{solution}{exo:b1:elementary-steps:3}
Het intermediair is het dal op 30; de \omterm{def:b1:elementary-steps:energy-profile}{overgangstoestanden} zijn de maxima
op 70 en 55. De hoogste \omterm{def:b1:elementary-steps:energy-profile}{overgangstoestand} is de eerste: de eerste stap
is snelheidsbepalend. Hij is endotherm (van 0 omhoog naar 30).
\end{solution}

\begin{solution}{exo:b1:elementary-steps:4}
Met dezelfde factor $10^4$ (\cref{prop:b1:elementary-steps:catalysis}); de
evenwichtsconstante blijft ongewijzigd; het evenwicht wordt ongeveer
$10^4$ keer sneller bereikt.
\end{solution}

\begin{solution}{exo:b1:elementary-steps:5}
$t_{\max} = \ln(0.30/0.10)/(0.30 - 0.10) = \qty{5.49}{s}$. Dan is
$[\mathrm{B}]_{\max} = a_0\frac{0.10}{0.20}(\eu^{-0.549} - \eu^{-1.648}) =
0.5\,a_0(0.577 - 0.192) = 0.192\,a_0$.
\end{solution}

\begin{solution}{exo:b1:elementary-steps:6}
$\dd[\mathrm{I}]/\dd t = k_1[\mathrm{A}] - (k_{-1} + k_2)[\mathrm{I}] = 0$
geeft $[\mathrm{I}] = k_1[\mathrm{A}]/(k_{-1} + k_2)$ en $v = k_2[\mathrm{I}]
= \frac{k_1k_2}{k_{-1}+k_2}[\mathrm{A}]$. Als $k_2 \gg k_{-1}$, is $v =
k_1[\mathrm{A}]$: de eerste stap is snelheidsbepalend. Als $k_2 \ll k_{-1}$,
is $v = (k_1/k_{-1})k_2[\mathrm{A}]$: een voorevenwicht gevolgd door een
trage stap.
\end{solution}

\begin{solution}{exo:b1:elementary-steps:7}
Som: \ce{2H2O2 -> 2H2O + O2}. \omterm{def:b1:elementary-steps:catalyst}{Katalysator}: \ce{I-} (verbruikt, daarna
teruggevormd); intermediair: \ce{IO-}. De trage eerste stap bepaalt de
snelheid: $v = k_1[\ce{H2O2}][\ce{I-}]$.
\end{solution}

\begin{solution}{exo:b1:elementary-steps:8}
De stap is sterk opwaarts: volgens het postulaat van Hammond lijkt zijn
\omterm{def:b1:elementary-steps:energy-profile}{overgangstoestand} op het carbokation. Een substituent die de energie van
het carbokation verlaagt, verlaagt die van de \omterm{def:b1:elementary-steps:energy-profile}{overgangstoestand} bijna
evenveel, zodat de barrière daalt en de stap sneller verloopt.
\end{solution}

\begin{solution}{exo:b1:elementary-steps:9}
$k_{\mathrm B}/k_{\mathrm C} = \exp(20000/RT)$: \num{3.0e3} bij \qty{300}{K},
$\exp(4.01) = 55$ bij \qty{600}{K}. Verwarmen vermindert de kinetische
voorkeur voor B en laat het systeem, doordat de vorming van B omkeerbaar
wordt, het stabielere C bereiken.
\end{solution}

\begin{solution}{exo:b1:elementary-steps:10}
$k_1[\mathrm{A}][\mathrm{M}] = (k_{-1}[\mathrm{M}] + k_2)[\mathrm{A}^*]$,
dus $v = k_2[\mathrm{A}^*] = \dfrac{k_1k_2[\mathrm{A}][\mathrm{M}]}{k_{-1}[\mathrm{M}]
+ k_2}$. Bij hoge druk ($k_{-1}[\mathrm{M}] \gg k_2$) is $v =
(k_1k_2/k_{-1})[\mathrm{A}]$: orde 1. Bij lage druk is $v =
k_1[\mathrm{A}][\mathrm{M}]$: orde 2, de activering door botsing wordt de
trage stap.
\end{solution}

\begin{solution}{exo:b1:elementary-steps:11}
Het bewijs is dat van de stelling. Voor $k_1 = k_2 = k$:
$\frac{\dd}{\dd t}([\mathrm{B}]\eu^{kt}) = ka_0$, dus $[\mathrm{B}]\eu^{kt} =
ka_0t$ en $[\mathrm{B}] = a_0kt\,\eu^{-kt}$, maximaal bij $t = 1/k$.
\end{solution}

\begin{solution}{exo:b1:elementary-steps:12}
$\mathrm{A} + \mathrm{B} \rightleftharpoons \mathrm{I} + \mathrm{C}$ (snel,
constante $K_1$), daarna $\mathrm{I} + \mathrm{B} \to \mathrm{P}$ (traag,
$k_2$). De som is de totale reactie; $[\mathrm{I}] =
K_1[\mathrm{A}][\mathrm{B}]/[\mathrm{C}]$ en $v = k_2[\mathrm{I}][\mathrm{B}]
= k_2K_1[\mathrm{A}][\mathrm{B}]^2/[\mathrm{C}]$: het bijproduct, dat in het
voorevenwicht vrijkomt, vertraagt de reactie.
\end{solution}

\begin{solution}{pb:b1:elementary-steps:1}
\textbf{1.} Een verdubbeling van $[\ce{NO}]_0$ vermenigvuldigt $v_0$ met
4: orde 2; een verdubbeling van
$[\ce{O2}]_0$ vermenigvuldigt ze met 2: orde 1.
\textbf{2.} $v = k[\ce{NO}]^2[\ce{O2}]$; $k = 1.0 \times 10^{-5}/((1.0
\times 10^{-3})^2 \times 1.0 \times 10^{-3}) =
\qty{1.0e4}{L^2.mol^{-2}.s^{-1}}$.
\textbf{3.} Een gelijktijdige botsing van drie moleculen is zeldzaam, en
één enkele stap zou bij verwarming niet trager kunnen worden (vraag 5).
\textbf{4.} $E_a = R\ln(k_{400}/k_{300})/(1/300 - 1/400) = 8.314 \ln(0.1)
/(8.33 \times 10^{-4}) = \qty{-23}{kJ/mol}$.
\textbf{5.} Een \omterm{def:b1:elementary-steps:elementary-step}{elementaire stap} verloopt via botsingen die een barrière
overschrijden; de fractie daarvan die dat kan, $\eu^{-E_a/RT}$ met $E_a \ge
0$, kan alleen met de temperatuur toenemen.
\textbf{6.} $-\dd[\ce{NO}]/\dd t = 2v$, $-\dd[\ce{O2}]/\dd t = v$.
\textbf{7.} \ce{2NO -> N2O2} plus \ce{N2O2 + O2 -> 2NO2} geeft
\ce{2NO + O2 -> 2NO2}; het intermediair is \ce{N2O2}.
\textbf{8.} Beide bimoleculair (en de omgekeerde van de eerste
unimoleculair).
\textbf{9.} $v = k_2[\ce{N2O2}][\ce{O2}]$.
\textbf{10.} $[\ce{N2O2}] = K_1[\ce{NO}]^2$ met $K_1 = k_1/k_{-1}$.
\textbf{11.} $v = k_2K_1[\ce{NO}]^2[\ce{O2}]$: de waargenomen wet.
\textbf{12.} $k = k_1k_2/k_{-1}$.
\textbf{13.} $\dd[\ce{N2O2}]/\dd t = k_1[\ce{NO}]^2 - k_{-1}[\ce{N2O2}] -
k_2[\ce{N2O2}][\ce{O2}]$.
\textbf{14.} Nul stellen geeft: $[\ce{N2O2}] = k_1[\ce{NO}]^2/(k_{-1} +
k_2[\ce{O2}])$.
\textbf{15.} $v = k_2[\ce{N2O2}][\ce{O2}]$ geeft de opgegeven wet.
\textbf{16.} $k_{-1} \gg k_2[\ce{O2}]$: het dimeer valt veel vaker uiteen
dan het dizuurstof ontmoet, zodat de eerste stap in evenwicht blijft.
\textbf{17.} Als $k_2[\ce{O2}] \gg k_{-1}$, is $v = k_1[\ce{NO}]^2$: orde
2 in \ce{NO} en 0 in \ce{O2}.
\textbf{18.} De gemeten orde 1 in \ce{O2}: het eerste grensgeval.
\textbf{19.} $\ln k = \ln k_1 + \ln k_2 - \ln k_{-1}$.
\textbf{20.} Differentiëren naar $T$ met $\dd\ln k_i/\dd T
= E_i/RT^2$ geeft $E_a = E_1 + E_2 - E_{-1}$.
\textbf{21.} Verwarmen versnelt de trage stap een beetje ($E_2$ klein),
maar verschuift het voorevenwicht sterk terug naar \ce{NO} ($E_{-1}$
groot): minder dimeren, en de netto snelheid daalt.
\textbf{22.} Het dimeer ligt $E_1 - E_{-1} = \qty{-38}{kJ/mol}$ onder de
beginstoffen en de tweede \omterm{def:b1:elementary-steps:energy-profile}{overgangstoestand} $E_2 = \qty{15}{kJ/mol}$
boven het dimeer: \qty{23}{kJ/mol} onder de beginstoffen. Het hoogste
punt van het pad is de eerste \omterm{def:b1:elementary-steps:energy-profile}{overgangstoestand}, maar \qty{2}{kJ/mol}
hoger.
\textbf{23.} $\exp\big(\frac{23000}{8.314}(\frac{1}{400} - \frac{1}{300})\big)
= 0.10$: tien keer trager bij \qty{400}{K}.
\textbf{24.} $E_a = 2 - 40 + 15 = \textbf{\qty{-23}{kJ/mol}}$, de waarde
die in Deel I gemeten werd: het mechanisme verklaart zowel de
\omterm{def:b1:rate-laws:order}{snelheidswet} als haar vreemde temperatuurafhankelijkheid.
\end{solution}
